\section{The non linear time focused operator}\label{sec_time_focus}

\subsection{Atoms and time focused transform}
\label{sse:definitions.timefocus}
Let us now introduce and study the second nonlinear transform of interest here, namely the time focused transform, which involves a focus function defined in the time domain. Even though the ideas behind the transform are the same as the ones developed for frequency focus, the definition and the method used to prove the needed results are completely different from the one used in Section~\ref{sec_frequency_focus}.


\paragraph{Assumptions}

One of the first differences with the frequency focused operator are the assumptions on the frequency function $\gamma$. In this section $\gamma$ is a $C^1$ symmetrical diffeomorphism satisfying $\lim_{t\to -\infty}\gamma(t)=-\infty$  and $\lim_{t\to +\infty}\gamma(t)=+\infty$. A prototypical $\gamma$ could be hyperbolic sinus shape like, $\ie$ very precise around $0$ (low frequencies) and less precise for high values (high frequencies).

We will assume the window $h$ to be continuous (at least piecewise) and with compact support of size $l\in\R^+_*$. An example of such window would be the Hamming window. About the focus function of $f$ - written $\sigma_f^\tau$ - we will assume it to be grater than $1$, continuous (at least piecewise) and only vanishing at infinity (not necessarily with compact support), \ie
\begin{equation}
    \forall f\in L^2(\R), \  \sigma^\tau_f \in C^+_{p,0}(\R) + 1.
\end{equation}
In order to lighten the notations we will sometimes omit the subscript $\tau$ when there are no ambiguities.
\begin{remark}
    Whereas in the previous transforms $Lf$ and $M^\nu$ the second variable $u$ is a 'frequency index', here $\omega$ is directly a frequency and the meaning of $\gamma$ is a bit different in this section. The function $\gamma$ is still homogeneous to a frequency, however it models a non linearity in the precision of the analysis along the frequency axis.
\end{remark}

\paragraph{Atoms and transform definition} Given hypotheses we can define the  time focused atoms .
\begin{equation}
	\forall x,t,\omega \in\R , \ h_{t,\omega,\sigma_f^\tau}(x):=  \sqrt{\gamma'(\omega)\sigma(t)} e^{2i\pi\gamma(\omega)x}h(\sigma_f^\tau(t)(x-t))\ ,
\end{equation}
and the transform of a given signal $f\in L^2$ by
\begin{equation}\label{eq_def_operator_time_focus}
	\forall t,\omega\in\R, \ M^\tau f(t,\omega) := \langle f , h_{t,\omega,\sigma_f^\tau} \rangle_{L^2_x}.
\end{equation}
The pointwise definition of the scalar product is guaranteed by the fact that $h(\sigma(t)(x-t))$ is a compactly supported function of $x$.

Sometimes, if there are no ambiguities we will write $h_{t,\omega,f}$ instead of $h_{t,\omega,\sigma_f}$ and $h_{t,\omega,n}$ instead of $h_{t,\omega,\sigma_{f_n}}$ for a certain sequence $(f_n)_n$.

\paragraph{Density result}

In order to justify as much as possible the validity of our further computations we will prove the following density result which state that if we take a $L^2$ function $f$ and a $C_c$ function $f_\eps$ as close as possible in $L^2$ to $f$ then the norm of $M^\tau f$ will be controlled by the norm of $M^\tau f_\eps$ which is finite thanks to basic integration rules.

\begin{lemma}\label{lemma_density_result}
    Let $f\in L^2$ and $(f_n)_n \in (C_c)^\N$ be such that $f_n \rightarrow f$ in $L^2$ and $\sigma(f_n)=:\sigma_n\rightarrow \sigma(f)$ in $L^\infty$ then
    \begin{equation}
        \forall \eps>0, \exists N\in\N, \forall n\geqslant N, \left|\|M^\tau f\|_{L^2} - \|M^\tau f_n \|_{L^2} \right| \leqslant \eps\ .
    \end{equation}
\end{lemma}

\begin{proof}
    We have
    \begin{align*}
        \left| \|M^\tau f_n \|^2_{L^2} - \|M^\tau f\|^2_{L^2} \right| & = \left| \iint_{\R^2} |\langle f_n , h_{t,\omega,n}\rangle|^2 d\omega dt - \iint_{\R^2} |\langle f , h_{t,\omega,f}\rangle|^2 d\omega dt \right| \\
        & = \left| \int_\R A(t)dt + \int_\R B(t)dt \right|,
    \end{align*}
    with
    \begin{align*}
        A(t) &:= \int_\R \sigma_f(t)h^2(\sigma_f(t)(x-t))\left(|f_n(x)|^2-|f(x)|^2\right) dx\ .
    \end{align*}
    Hence
    \begin{equation*}
        \int_\R |A(t)|dt \leqslant \|\sigma_f\|_\infty \|h(\sigma_f(t)(x-t))\|_{L^\infty_{x,t}}\eps\ .
    \end{equation*}
    Similarly 
    \begin{equation*}
    B(t):= \int_\R \sigma_n(t)|f_n(x)|^2\left(h^2(\sigma_n(t)(x-t))-h^2(\sigma_f(t)(x-t))\right)dx\ ,
    \end{equation*}
    thus
    \begin{equation*}
        \int_\R |B(t)|dt \leqslant (\|\sigma_f\|_\infty+\eps)(\|f\|_{L^2}+\eps)^2 2l\eps\ .
    \end{equation*}
\end{proof}

The result alone does not allow us to extend directly the well-definedness from $C_c$ to $L^2$. However we will see that combined with the upper bound control proved in Proposition~\ref{prop_norm_control_time_operator_L_1_lemma} it will allow us to conclude about the well-definedness of $M^\tau$ on $L^2$.

\subsection{Example and ideas of the time focus function}\label{ssec_time_focus_function}
Again, it is not the goal of the current paper to discuss in details explicit choices for the focus functions that would be relevant in specific applications. We only provide a simple and generic example that could be of interest (a more sophisticated alternative, developed in the context of audio compression, can be found in~\cite{Molla2004determining}).
\begin{example}[Adaptation to transient events]
\label{ex:time_focus}
As already mentioned in the introduction, a fairly natural idea is to increase the time resolution of the analysis when the analysed signal has faster variations. This requires a heuristics to measure such speed of variations. As in Example~\ref{ex:frequency_focus}, a ratio of suitably chosen norms can be used for that purpose. Using frequency content as a measure of speed of variations, one may consider the quantity 
\begin{equation}\label{eq_def2_sigma_tau}
	\sigma_{f}^\tau(t) := 1 + \sqrt{\frac{\int_\R |\omega|^{2n} |V_g f(t,\omega)|^2 d\omega}{\int_\R |V_g f(t,\omega)|^2d\omega}}\varphi(t)\ ,
\end{equation}
with $V_gf$ the STFT of $f$ with gaussian window $g(x):=e^{-x^2/2a^2}$ with variance $a^2$, and $\varphi$ an additional function, which can ensure constant support and appropriate normalisation. The rationale is to use the ratio to measure some frequency content of the signal around time $t$ (as an average frequency variable, weighted by the spectrogram at time $t$). The higher the average frequency, the larger the time focus, and smaller the adapted window. A toy example is given in Appendix~\ref{sec:app.time.foc.example} that shows the ability of such functions to detect transient events in a signal and therefore adapt the scale of a window in the time focused transform.

\end{example}

\subsection{Norm relationship and time kernel}

Following the same ideas as in Section~\ref{sec_frequency_focus} we will determine an equality of the $L^2$ norm of $Mf$ and introduce a certain time dependent kernel. Beside being useful for the rest of the study the equality guarantees the well-definedness of the transform on $L^2$.

However one main difference with the frequency kernel is that we will be able to have an explicit expression or at least a way easier to manipulate expression of the kernel. The first result gives us the norm relationship in Fourier.

\begin{prop}
    Using the previous definition, for $f\in C_c$ we have then
	\begin{equation}
		\|M^\tau f\|_{L^2}^2 = \iint_{\R^2} \hat{f}(\xi)\bar{\hat{f}}(\xi') K_{\sigma_f^\tau}(\xi-\xi') d\xi d\xi',
	\end{equation}
	where the time kernel\footnote{Even if we call it the time-kernel, its variable is still frequency since we will take $\F^{-1}$. This has to be taken as a abuse of language} is 
	\begin{equation}\label{eq_def_time_kernel}
		K_{\sigma_f^\tau}(u) := \iint_{\R^2}e^{-2i\pi u t}\bar{\hat{h}}(z)\hat{h}\left(z-\tfrac{u}{\sigma_f^\tau(t)}\right)\, dzdt\ .
	\end{equation}
\end{prop}
\begin{proof}
	In order to make notations lighter we use $\sigma$ for $\sigma_f^\tau$. The computations are quite similar than the one used in the proofs of the frequency focus.
	\begin{align*}
		\|M^\tau f\|_{L^2}^2 & = \iint_{\R^2} |\langle \hat{f},\hat{h}_{t,\omega,\sigma}\rangle|^2dtd\omega \\
		& = \iiiint_{\R^4} \hat{f}(\xi)\bar{\hat{f}}(\xi')\bar{\hat{h}}_{t,\omega,\sigma}(\xi)\hat{h}_{t,\omega,\sigma}(\xi') d\xi d\xi' dt d\omega \\
		& = \iint_{\R^2}\hat{f}(\xi)\bar{\hat{f}}(\xi') \underset{=:K_{\sigma}(\xi,\xi')}{\underbrace{\iint_{\R^2}\bar{\hat{h}}_{t,\omega,\sigma}(\xi)\hat{h}_{t,\omega,\sigma}(\xi')dt d\omega}} d\xi d\xi'
	\end{align*}
	We have
	\begin{equation*}
		K_{\sigma}(\xi,\xi'):= \iint_{\R^2}\frac{\gamma'(\omega)}{\sigma(t)}e^{-2i\pi(\xi-\xi')t}\bar{\hat{h}}\left(\tfrac{\xi-\gamma(\omega)}{\sigma(t)}\right)\hat{h}\left(\tfrac{\xi'-\gamma(\omega)}{\sigma(t)}\right)\,d\omega dt\ .
	\end{equation*}
	Hence, by setting $y(\omega)=-\frac{\gamma(\omega)}{\sigma(t)}$ we obtain
	\begin{align*}
		K_{\sigma}(\xi,\xi') &= \iint_{\R^2} \bar{\hat{h}}\left(\tfrac{\xi}{\sigma(t)}+y\right)\hat{h}\left(\tfrac{\xi'}{\sigma(t)}+y\right)\,e^{-2i\pi(\xi-\xi')t}\,dydt \\
		&= \iint_{\R^2}e^{-2i\pi(\xi-\xi')t} \bar{\hat{h}}(z)\hat{h}\left(z-\tfrac{\xi-\xi'}{\sigma(t)}\right)\,dzdt\ .
	\end{align*}
	The last equality holds by applying the changing of variable $z=y+\frac{\xi}{\sigma(t)}$. And now since the kernel is a function of $\xi-\xi'$ we can deduce the expected result.
\end{proof}

\begin{corollary}
	We have for $f\in C_c(\R)$,
	\begin{equation}
		\|M^\tau f\|_{L^2}^2 = \int_\R |f|^2(t)\F^{-1}(K_{\sigma_f^\tau})(t) dt.
	\end{equation}
\end{corollary}
\begin{proof}
	The proof use the Plancherel equality and the fact that the Fourier transform of a convolution product is the product of the Fourier transform. It is an elementary computation
	\begin{align*}
		\|M^\tau f\|_{L^2}^2  &= \iint_{\R^2} \hat{f}(\xi)\bar{\hat{f}}(\xi') K_{\sigma}(\xi-\xi') d\xi d\xi \\
		& = \int_\R \hat{f}(\xi) \int_\R \bar{\hat{f}}(\xi')K_\sigma(\xi-\xi')d\xi' d\xi \\
		& = \langle \hat{f} , K_\sigma \ast \bar{\hat{f}} \rangle_{L^2_\xi} \\
		& = \int_\R f(t) \F^{-1}\left( K_\sigma \ast \bar{\hat{f}} \right)(t) dt \\
		& = \int_\R |f|^2(t)\F^{-1}(K_\sigma)(t)dt. 
	\end{align*}
\end{proof}

Now we can see that the problem can be solved by controlling the kernel $K_\sigma$. In order to obtain such control we can introduce the explicit formula of $\F^{-1}(K_\sigma)$.
\begin{equation}\label{eq_formula_inverse_fourier_kernel}
		\forall t\in\R, \ \F^{-1}(K_\sigma)(t) = \int_\R \sigma(x)|h(\sigma(x)(x-t))|^2dx.
\end{equation}


\subsection{Main result : norm control}

The rest of the section is dedicated to the proof of Theorem~\ref{theo_bounding_time_marseillan_transform} below. This result is the analogous control that we have in section~\ref{sec_frequency_focus} with  Theorem~\ref{theorem_norm_control_frequency_focus}.


\begin{theorem}\label{theo_bounding_time_marseillan_transform}
	Let $f\in L^2(\R)$ and $\sigma_f^\tau$ be a time focus parameter, then
	\begin{equation}
		c_f \|f\|_{L^2} \leqslant\| M^\tau f\|_{L^2_{t,\omega}} \leqslant C_f \|f\|_{L^2}\ ,
	\end{equation} 
	where
	\begin{equation}
		c_f^2 = \inf_{t\in\R}\left\{ \int_{-m}^{m}|h(x\sigma_f^\tau(x+t))|^2 dx \right\}>0 \ ; \ m:=\min\{1/\|\sigma_f^\tau\|_\infty,1\}
	\end{equation}
	and
	\begin{equation*}
		C_f^2 = \int_\R \frac{|h(\sigma_f^\tau(t)t)|^2}{\sigma_f^\tau(t)}dt < \infty
	\end{equation*}
\end{theorem}

This result is the same kind than the one obtained in the frequency focus situation. And still interesting, even though the control depends on the signal the dependence is made through the focus parameter $\sigma_f^\nu$. Hence it leads to the ideas of further studies about the control of this parameter which could help us to inverse such non linear operator.

We will see that the proof of such theorem \ref{theo_bounding_time_marseillan_transform} is way more direct than the frequency one. In the sens that we do not need to work on approximation of $\sigma$ by steps function.


\subsection{Upper bound control}

In order to prove the upper bound control - which by the way guarantees the well-definedness in $L^2(\R)\rightarrow L^2(\R^2)$ of our operator - we will control the $L^1$ norm of the previously introduced kernel $K_\sigma$. Until we extend the result to $L^2$, $f$ denotes a $C_c$ function.

\begin{lemma}[$L^1$ norm control]\label{lemma_norm_control_time_kernel_L1}
	Let $\sigma(t)$ be a focus function and $h$ be a window function, as defined in Section~\ref{sse:definitions.timefocus}. Let the time kernel $K_\sigma(u)$ be defined as in equation~\eqref{eq_def_time_kernel}.
%	\begin{equation}
%		K_\sigma(u) = \iint_{\R^2} e^{-2i\pi ut}\bar{\hat{h}}(z)\hat{h}(z-\tfrac{u}%{\sigma(t)})dzdt.
%	\end{equation}
	Then
	\begin{equation}
		\int_\R K_\sigma(u)du = \int_{\R} |h(\sigma(t)t)|^2 \frac{dt}{\sigma(t)}\ .
	\end{equation}
\end{lemma}
\begin{proof}
	We have by Fubini's theorem and setting $u' = z - \tfrac{u}{\sigma(t)}$,
	\begin{align*}
		\int_\R K_\sigma(u)du & = \iint_{\R^2} \bar{\hat{h}}(z)\int_{\R} e^{-2i\pi ut}\hat{h}(z-\tfrac{u}{\sigma(t)}) du dzdt \\
		& = \iint_{\R^2} \bar{\hat{h}}(z) \int_{\R} e^{-2i\pi \sigma(t)t(z-u)}\hat{h}(u) du\frac{dt}{\sigma(t)}dz \\
		& = \iint_{\R^2} \bar{\hat{h}}(z)e^{-2i\pi \sigma(t)tz}\int_{\R}\hat{h}(u)e^{2i\pi \sigma(t)tu}dudz\frac{dt}{\sigma(t)}\\
		& = \iint_{\R^2} \bar{\hat{h}}(z)e^{-2i\pi \sigma(t)tz}h(\sigma(t)t) dz\frac{dt}{\sigma(t)} \\
		& = \int_{\R} |h(\sigma(t)t)|^2\frac{dt}{\sigma(t)}
	\end{align*}
\end{proof}

From Lemma~\ref{lemma_norm_control_time_kernel_L1} we obtain the upper norm control by introducing a weighed window $H_f$ which lightens a bit the notations.
\begin{prop}[Upper bound]\label{prop_norm_control_time_operator_L_1_lemma}
	Let $f\in L^2(\R)$ and set the weighed window
	\begin{equation}
		H_f(t) := \frac{|h(\sigma_f(t)t)|}{\sqrt{\sigma_f(t)}},
	\end{equation}
	then we have
	\begin{equation}
		\|Mf\|_{L^2} \leqslant \|f\|_{L^2}\|H_f\|_{L^2}.
	\end{equation}
\end{prop}
\begin{proof}For the moment let us assume $f\in C_c$. For any $f,g\in L^2$ the dual equality
	\begin{equation*}
		\int_{\R} f(t)\hat{g}(t)dt = \int_{\R} \hat{f}(\omega)g(\omega)d\omega
	\end{equation*}
	gives
	\begin{equation*}
		\|Mf\|_{L^2}^2 = \int_{\R} \F(|f|^2)(\omega)K_\sigma(\omega)d\omega\ .
	\end{equation*}
	As $f\in L^2$ we have $\F(|f|^2)\in L^\infty$, hence by using Lemma~\ref{lemma_norm_control_time_kernel_L1} for the second member
	\begin{align*}
		\int_{\R} \F(|f|^2)(\omega)K_\sigma(\omega)d\omega & \leqslant \| \F(|f|^2)\|_{L^\infty}\int_{\R} |h(\sigma(t)t)|^2\frac{dt}{\sigma(t)}, \\
		& \leqslant \| |f|^2 \|_{L^1} \| H_f\|_{L^2}^2 \\
		& = \|f\|_{L^2}^2 \|H_f \|_{L^2}^2.				
	\end{align*}
	By the definition of $H_f$ and taking its $L^2$ norm we obtain the bound in Proposition~\ref{prop_norm_control_time_operator_L_1_lemma}. Therefore the upper bound $C_f$ in Theorem~\ref{theo_bounding_time_marseillan_transform} equals $\|H_f\|_{L^2}$.

 We will now extend the result to $L^2$. Let $f\in L^2$ and introduce a sequence $(f_n)_n$ with $f_n\in C_c$, converging to $f$ in $L^2$ and such that $\sigma(f_n)$ converges to $\sigma(f)$ in $L^\infty$. We know by Lemma~\ref{lemma_density_result} that 
 \begin{equation}
     \left|\|M^\tau f\|_{L^2} - \|M^\tau f_n \|_{L^2} \right| \rightarrow 0
 \end{equation}
 However we know by the upper bound control that
 \begin{equation}
     \|M^\tau f_n\|_{L^2} \leqslant C_n\|f_n\|_{L^2} \leqslant \sup_n \{C_n\}\|f_n\|_{L^2}
 \end{equation}
 Since $\sigma(f_n)$ converges in $L^\infty$ we know that $\sup_n \{C_n\}< \infty$. Hence $(\|M^\tau f_n\|)_n$ is bounded, and Lemma~\ref{lemma_density_result} guarantees that the $L^2$ norm of $M^\tau f$ for $f\in L^2$ is finite and controlled by $\|f\|_{L^2}C_f$.
 
\end{proof}

The homogeneity is respected in those results, indeed $\|Mf\|_{L^2}^2$ depends quadratically on $h$, \ie\, for any $\lambda\in\C$ we have $\|M_{\lambda h}f\|_{L^2}^2 = \lambda^2 \|M_h f\|_{L^2}^2$, so that $M_{\lambda h}f$ also satisfies Proposition~\ref{prop_norm_control_time_operator_L_1_lemma}.



The following norm control is not necessary at this point but may be useful in future work so it is presented here.
\begin{prop}[$L^2$ norm control]
	We keep the same hypothesis and notations of the time kernel \ref{eq_def_time_kernel}, thus we have
	\begin{equation}
		\|K_\sigma \|_{L^2}^2 = \|h\|_{L^2}^2 \int_{\R} |h(\sigma(t)t)|^2\frac{dt}{\sigma^2(t)}.
	\end{equation}
\end{prop}
\begin{proof}
	Using Fubini's lemma we obtain
	\begin{equation*}
		\|K_\sigma\|_{L^2}^2 = \iiiint_{\R^4} e^{-2i\pi(u-u')t}|\hat{h}(z)|^2\hat{h}(z-\tfrac{u}{\sigma(t)})\bar{\hat{h}}(z-\tfrac{u'}{\sigma(t)})dudu'dtdz
	\end{equation*}
	We set then $x:=\tfrac{u}{\sigma(t)}$ and $x'=\tfrac{u'}{\sigma(t)}$ which gives
	\begin{align*}
		\|K_\sigma\|_{L^2}^2  &= \iint_{\R^2} \underset{(1)}{\underbrace{\left( \int_{\R} e^{-2i\pi x\sigma(t)t}\hat{h}(z-x)dx \right)}}\underset{(2)}{\underbrace{\left( \int_{\R} e^{2i\pi x'\sigma(t)t}\bar{\hat{h}}(z-x')dx' \right)}}|\hat{h}(z)|^2 \frac{dt}{\sigma^2(t)}dz.
	\end{align*}
	We set $x_2=z-x$ and hence we recognise the inverse Fourier transform formula
	\begin{align*}
		(1) &= \int_{\R} e^{-2i\pi(z-x_2)\sigma(t)t}\hat{h}(x_2)dx_2 \\
		& = e^{-2i\pi z\sigma(t)t}h(\sigma(t)t).
	\end{align*}
	Furthermore $(2) = \overline{(1)} $. Thus when we gather all the members we obtain by Fourier isometry
	\begin{align*}
		\|K_\sigma\|_{L^2}^2 &= \iint_{\R^2} |h(\sigma(t)t)|^2 |\hat{h}(z)|^2 dz \frac{dt}{\sigma^2(t)} \\
		& = \|h\|_{L^2}^2 \int_{\R} |h(\sigma(t)t)|^2\frac{dt}{\sigma^2(t)}.
	\end{align*}
\end{proof}




\subsection{Lower bound control}

The following proposition gives a strictly positive lower bound for the norm, in the case of non zero signals.
\begin{prop}[Lower bound]\label{prop_lower_bound_control_time_variance}
	Let $\sigma$ be as previously. There exists $c=c(\sigma)>0$ which depends on $\sigma$ such that
	\begin{equation}\label{eq_lower_bound_kernel}
		\forall t\in\R, \ \F^{-1}(K_\sigma)(t) > c\ ,
	\end{equation}
	with
	\begin{equation}
		c:= \inf_{t\in\R}\left\{ \int_{-m}^{m}|h(x\sigma(x+t))|^2 dx \right\}\ ,
	\end{equation}
	where $m:=l/\|\sigma\|_{\infty}$ and $l\in\bar{\R}$ being such that $\text{supp}(h)\subset [-l,l]$.
\end{prop}
\begin{proof}
    Let us introduce the weighted norm function
\begin{equation}
    \forall t\in\R, \ H(t) := \int_{-m}^{m} |h(x\sigma(x+t))|^2 dx\ .
\end{equation}
%
Let $t\in\R$. Since $\sigma(x)\geqslant1$ for any $x\in\R$, by using equation~\eqref{eq_formula_inverse_fourier_kernel} we have
\begin{align*}
    \mathcal{F}^{-1}(K_\sigma)(t) & \geqslant \int_\R |h(\sigma(x)(x-t))|^2 dx \\
    & = \int_\R |h(\sigma(x+t)x)|^2 dx\ .
\end{align*}
Furthermore, if $|x|\geqslant l/\|\sigma\|_\infty$ then $h(x\sigma(x+t))=0$ and thus
\begin{equation}
    \int_\R |h(\sigma(x+t)x)|^2 dx = \int_{-m}^m |h(\sigma(x+t)x)|^2 dx\ .
\end{equation}
%
Hence we have
\begin{equation}
    \forall t\in\R, \ \mathcal{F}^{-1}(K_\sigma)(t)\geqslant H(t).
\end{equation}
Now let us prove that $\inf_t H(t) =: c >0$. We obviously have $H(t)>0$ for any $t\in\R$ furthermore by the use of the dominated convergence theorem we also have $\lim_{\pm\infty}H(t)>0$. And since $H(t)$ is continuous (still the use of the dominated convergence theorem) we can conclude that there exists $c=c(\sigma)>0$ such that
\begin{equation}
    \inf_{t\in\R}H(t)=c>0\ .
\end{equation}
\end{proof}

